\textbf{Pour cette première partie, aucune justification n'est demandée et une seule réponse est possible par question. Pour chaque question, reportez son numéro sur votre copie et indiquez votre réponse.}

	Une réponse fausse ou l'absence de réponse n'enlève aucun point.


	\vspace{0,25cm}
\subsubsection*{Question 1}
\mbox{$\dfrac25-\dfrac{3}{10}=$}
		\begin{enumerate}[label=\alph*)]
\begin{multicols}{4}
			\item \mbox{$-\dfrac{3}{25}$}
			\item \mbox{$-\dfrac15$}
			\item \mbox{$\dfrac{1}{10}$}
			\item \mbox{$\dfrac15$}
				\end{multicols}
		\end{enumerate}


	\textcolor{red}{\textsc{Correction}}\\

	\mbox{$\dfrac25-\dfrac{3}{10}=\dfrac{4}{10}-\dfrac{3}{10}=\dfrac{1}{10}$}, \framebox[1.2\width]{réponse C}.
		\subsubsection*{Question 2}

Laquelle de ces fonctions est représentée graphiquement par une droite ?		\begin{enumerate}[label=\alph*)]
			\begin{multicols}{4}
				\item \mbox{$f(x)=\dfrac52 x-5$}
				\item \mbox{$g(x)=x^3$}
				\item \mbox{$h(x)=-\dfrac1 x +3$}
				\item \mbox{$i(x)=2x^2+3x+1$}

				\end{multicols}
		\end{enumerate}


\textcolor{red}{\textsc{Correction}}\\
Une fonction \underline{affine} a pour représentation graphique une droite : il s'agit donc de \mbox{$f(x)=\dfrac52 x-5$}, \framebox[1.2\width]{réponse A}.
\subsubsection*{Question 3}
\begin{minipage}[t]{0.45\linewidth}\vspace{0mm}
	Dans le repère ci-contre sont tracées les droites $(d_1)$, $(d_2)$, $(d_3)$ et $(d_4)$.

	La droite ayant pour équation réduite \mbox{$y=-\dfrac12 x +1$} est :
	\begin{enumerate}[label=\alph*)]
		\begin{multicols}{2}
			\item $(d_1)$
			\item $(d_2)$

		\end{multicols}
		\begin{multicols}{2}
			\item $(d_3)$
			\item $(d_4)$
		\end{multicols}
	\end{enumerate}

\end{minipage}\hspace{0.05\linewidth}\begin{minipage}[t]{0.45\linewidth}\vspace{0mm}

	\begin{tikzpicture}[xscale=0.5,yscale=0.5,rotate=0,every node/.style={scale=0.6,rotate=0}]
		\clip(-6.4,-3.8)rectangle(10.4,5.4);
		\draw[blue!30,thick](-6.4,-3.8)grid(10.4,5.4);
		\draw[blue, line width=0.3mm,->, >=triangle 45](-6.4,0)--(10.4,0);
		\draw[blue, line width=0.3mm,->, >=triangle 45](0,-3.8)--(0,5.4);
		\foreach \i in {-6,-5,...,-1,1,2,...,10}
		{\draw(\i,-0.2)node{\textbf{\i}}(-0.2,\i)node{{$\i$}};}
		\draw(-0.2,-0.2)node{{$0$}};
		\draw[red, line width=0.3mm](-6.4,-2.2)--(10.4,6.2)node[pos=0.83, below]{$(d_1)$};
				\draw[green, line width=0.3mm](-6.4,-3.7)--(10.4,4.7)node[pos=0.83, below]{$(d_2)$};
				\draw[purple, line width=0.3mm](-6.4,-0.5)--(10.4,-0.5)node[pos=0.83, below]{$(d_3)$};
						\draw[olive, line width=0.3mm](-6.4,4.2)--(10.4,-4.2)node[pos=0.83, above]{$(d_4)$};
	\end{tikzpicture}

\end{minipage}



\textcolor{red}{\textsc{Correction}}\\
La droite d'équation \mbox{$y=-\dfrac12 x +1$} a une pente négative (donc elle \og{}descend\fg{}) : il s'agit donc nécessairement de $(d_4)$, \framebox[1.2\width]{réponse D}.

\subsubsection*{Question 4}
Un automobiliste roule à 60~km/h de moyenne pendant 2~h~30~min. La distance parcourue est :
\begin{enumerate}[label=\alph*)]
	\begin{multicols}{4}
		\item 150~km
		\item 140~km
		\item 130~km
		\item 120~km
				\end{multicols}
\end{enumerate}


\textcolor{red}{\textsc{Correction}}\\
2~h~30~min, c'est deux heures et demi, donc la distance parcourue est \mbox{$60\times\np{2.5}=150$}~km, \framebox[1.2\width]{réponse A}.

\subsubsection*{Question 5}
Un article à 200~\euro{} coûtera après une augmentation de 20\,\% :
\begin{enumerate}[label=\alph*)]
	\begin{multicols}{4}
		\null \vfill\item \mbox{$200+\np{0.2}$} \vfill \null
		\columnbreak
		\null \vfill\item \mbox{$200\times \dfrac{20}{100}$} \vfill \null
		\columnbreak
		\null \vfill\item \mbox{$200\times\np{1.2}$} \vfill \null
		\columnbreak
		\null \vfill\item \mbox{$100\times\np{0.20}+200$} \vfill \null

	\end{multicols}
\end{enumerate}


\textcolor{red}{\textsc{Correction}}\\
Augmenter une valeur de 20\,\% revient à multiplier celle-ci par \mbox{$\dfrac{100+20}{100}=\np{1.2}$}, \framebox[1.2\width]{réponse C}.
\subsubsection*{Question 6}
On développe \mbox{$(2x-5)^2$}. L'expression développée et réduite est :		\begin{enumerate}[label=\alph*)]
	\begin{multicols}{4}
		\item \mbox{$2x^2-10+25$}
		\item \mbox{$4x^2-20x-25$}
		\item \mbox{$4x^2-20x+25$}
		\item \mbox{$4x^2-25$}
		\end{multicols}
\end{enumerate}


\textcolor{red}{\textsc{Correction}}\\
C'est une identité remarquable, \mbox{$(2x-5)^2=(2x)^2-2\times2x\times5+5^2=4x^2-20x+25$},\framebox[1.2\width]{ réponse C}.
\subsubsection*{Question 7}

	La résistance $R$ (en ohms) d'un appareil est donnée par la formule \mbox{$R=\dfrac{U^2}{P}$} où $U$ est la tension (en volts) et $P$ la puissance (en watts).

	Quelle est la résistance d'un appareil lorsque la tension est de 20~volts et la puissance de 80~watts?

		\begin{enumerate}[label=\alph*)]
		\begin{multicols}{4}
		\null\vfill	\item $\dfrac12$ \vfill \null
			\columnbreak
		\null\vfill		\item $2$ \vfill \null
			\columnbreak
		\null\vfill		\item $\dfrac15$ \vfill \null
			\columnbreak
		\null\vfill		\item $5$\vfill\null

		\end{multicols}
	\end{enumerate}


	\textcolor{red}{\textsc{Correction}}\\
	On calcule \mbox{$R=\dfrac{20^2}{80}=\dfrac{20\times{20}}{4\times{20}}=\dfrac{20}{4}=5$}, \framebox[1.2\width]{réponse D}.
\subsubsection*{Question 8}
On considère les nombres suivants : \quad \mbox{$A=\dfrac18$},\quad \mbox{$B=\dfrac19$},\quad \mbox{$C=\dfrac{1}{12}$}\quad et \quad \mbox{D$=\np{0.1}$}.

Le plus petit de ces nombres est :
\begin{enumerate}[label=\alph*)]
	\begin{multicols}{4}
		\item A
		\item B
		\item C
		\item D
		\end{multicols}
\end{enumerate}


\textcolor{red}{\textsc{Correction}}\\
Puisque \mbox{$D=\dfrac{1}{10}$}, on doit comparer quatre fractions ayant le même numérateur $1$, la plus petite est celle qui a le plus grand dénominateur \mbox{$\dfrac{1}{12}$}, \framebox[1.2\width]{réponse C}.
\vspace{10mm}
\begin{center}
\end{center}

\vspace{10mm}
